

\label{98-Conclusioni}
Il progetto, una volta finito, è in grado di rivelare la maggior parte degli schiocchi effettuati (71\%): problemi sorgono tuttavia quando non si effettua uno schiocco, infatti muovendo leggermente la mano viene effettuata talvolta una lettura sbagliata, soprattutto quando viene mosso il pollice (sul quale è presente un sensore). In questo caso, il movimento non viene classificato come un anomalia ma come uno schiocco, causando l'accensione dei led.
Alcune fra le possibili cause alla base di questa problematica potrebbero essere:\\

\noindent{}Il modello utilizzato è in una condizione di "overfitting", ovvero il modello si è abituato ai dati dati nel training set, e classifica in modo errato i dati effettivi; un segnale d'allarme per questo particolare caso è dato dal "Model Testing" su Edge Impulse che riporta un'accuratezza \textbf{teorica} del modello pari al 80\%.\\

\noindent{}I dati sulle anomalie sono insufficienti. Semplicemente raccogliendo più dati, sopratutto su molti casi nei quali l'utente \textbf{non} schiocca le dita, è possibile allenare il modello a distinguere le due classi in maniera più accurata.\\

\noindent{}I sensori utilizzati non sono ottimali per il lavoro da eseguire; i sensori di Hall riportano una grande variazione ad un piccolo movimento se il sensore è molto vicino alla calamita (< 5 cm). Diversamente, ad una distanza maggiore di 5 cm il sensore rivelerà una piccolissima variazione di tensione in caso di movimento del magnete (al minimo 1.3 mV), mentre ad una distanza superiore ai 7.5 cm non verrà rilevata nessuna variazione di tensione.\\

\noindent{}Infine, i dati non stati modificati in alcun modo prima di essere passati alla rete neurale, quindi creando un nuovo "Processing Block" su Edge Impulse, oppure modificando direttamente i dati inviati dall'Arduino, è possibile facilitare l'addestramento della rete neurale.\\